The increasing integration of digital technologies into medical devices has introduced new challenges for quality assurance. To address these challenges, \gls*{ai}-based defect detection has gained attention as an efficient alternative to conventional rule-based and manual inspection methods. In the context of software reliability and cybersecurity, \cite{kim2024deep} proposed an AI-based software defect detection method as part of a broader framework for strengthening medical device security. Their work highlights the potential of \gls*{ai} to provide a more efficient alternative to conventional rule-based approaches for identifying software flaws that may compromise the reliability and safety of medical devices. From a manufacturing perspective, \cite{bonatti2025expert} developed an automated expert system for detecting porosity defects in prosthetic components produced using electron beam melting. By analyzing layer-by-layer optical images acquired during manufacturing, their approach enables rapid and nondestructive identification of defects and demonstrates strong agreement with computed tomography-based inspection. More recently, \cite{huang2026unsuperviseddefectdetectionsurgical} investigated unsupervised defect detection for surgical instruments and identified several challenges associated with applying conventional industrial defect detection methods to the medical device domain, including textured backgrounds, small defects, and domain shifts. Their proposed approach combines background masking, patch-based analysis, and domain adaptation to improve the detection of fine-grained defects. Collectively, these studies demonstrate the growing role of AI-based methods in improving the efficiency of defect detection across both the software and physical components of medical devices, while emphasizing the importance of reliable detection mechanisms in safety-critical healthcare applications.

As many modern medical devices rely on semiconductor components, advances in industrial semiconductor defect detection are also relevant to the medical device domain. Neural methods have been widely applied to defect detection, fault classification, and anomaly prediction in semiconductor manufacturing, particularly for analyzing process and inspection data. Prior work \cite{lee2017convolutional,khader2018stencil,lu2019realtime,cao2024xscan} has focused on improving detection accuracy and enabling real-time process control, including for solder paste printing, a major source of downstream defects. More recently, \cite{shenoy2026uncertainty} introduced a framework for generating structured causal explanations for root cause analysis, incorporating uncertainty to quantify the risk associated with inferred defect causes.

Crucially, these works do not explicitly account for uncertainty arising from the model's decisions themselves, i.e., epistemic uncertainty, which constitutes a central focus of our work.


% Rejection option has been studied since the dawn of artificial intelligence (cf.~\cite{Flores1958TEC, Chow1970TIT}) as a way to avoid low-confidence decisions~\cite{PudilNBK1992ICPR}, and miclassifications.

% Some methods require dedicated architectures \cite{CorbiereTBCP2019NeurIPS,GeifmanE2019ICML}.

% Some dedicated training data \cite{HuangZZ2020NeurIPS} proposed a method to improve the generalization empirical risk minimization of deep models, focusing on the task of learning from corrupted data and showing how calibrating the model at training time improves the models' performance for selective classification. A similar results is also reaffirmed in~\cite{FischJB2022CORR}.
% Recently, \cite{FengAHA2022CORR} have noticed that complex training techniques involving ad-hoc loss functions do not necessarily imply dramatic improvements for decision making with rejection option.

% methods like \cite{GraneseRGPP2021NeurIPS}

% strategies like \cite{markowitz1952portfolio,liu2019deep} we add a detailed analysis.

% \cite{protookn,oknregistry} \cite{fda_510k,fda_maude}
